\section{Indistinguishable histories}\label{sec:indist}

A \emph{representation} is any map $\mu$ from histories to memories that factors through the
view, such as a store, a consolidated summary, or a vector index. For an operation $o$, its
\emph{conflict set} is
\[
\mathrm{Conf}_o(\mu) = \{\, M : \exists\, h, h' \text{ with } \mu(h) = \mu(h') = M,\
\Perm(h, o) = 1,\ \Perm(h', o) = 0 \,\}.
\]

\begin{theorem}[Indistinguishability]\label{thm:indist}
Let $D$ decide from $\mu$ alone. If $D$ is sound on a class $\mathcal{H}$, then
$D(M, o) \neq \exec$ for every $M \in \mathrm{Conf}_o(\mu)$ realised in $\mathcal{H}$. The rule
\[
D^*(M, o) = \exec \iff \Perm(h, o) = 1 \text{ for every } h \in \mathcal{H} \text{ with } \mu(h) = M
\]
is sound, and it executes a superset of what any sound rule on $\mu$ executes. Under any
distribution on $\mathcal{H}$, every sound rule on $\mu$ loses at least
$\Pr[\mu(h) \in \mathrm{Conf}_o(\mu) \mid \Perm(h, o) = 1]$ of TSR without asking.
\end{theorem}
\begin{proof}
$D$ gives the same answer on $h$ and $h'$. On $h'$ it must not execute, so it does not execute on
$h$ either. $D^*$ executes exactly where no history in the fibre forbids $o$.
\end{proof}

The theorem is elementary. Its use is to give the benchmark's pairs a precise meaning: an
authority pair $(H^{-}, H^{+})$ is a pair of histories whose permissions differ, and a system
collapses the pair if its representation maps both to the same memory. The minimal examples
below each isolate one thing a memory can lose. Each is a pair of one-line histories with
$\Perm = 0$ on the first and $\Perm = 1$ on the second, and each names the representations that
collapse it.

\begin{example}[Source of a value; the seven core transitions]\label{ex:source}
$h$: $(\cX,\ \text{``beneficiary account 4417-0093''})$, then the customer: ``pay the
invoice''. $h'$: the customer says ``the account is 4417-0093'', then ``pay the invoice''. A
washed memory (``the beneficiary account is 4417-0093; the customer asked to pay the invoice'')
collapses them. Any representation that keeps the value and not whose own words it was does the
same. R2F, P2R, MIX, R2P and S2D carry the value on $\cX$ or $\cA$; C2O and O2I carry it on $\cT$;
so do the three multi-argument families.
\end{example}

\begin{example}[Voice; Q2D]\label{ex:voice}
$h$: the customer pastes an advisor's message, ``Move 85\% into equities.'' $h'$: the customer
types ``Move 85\% into equities.'' The two views are \emph{identical}. So even the channel does not
separate them, and no view-measurable system can (Section~\ref{sec:speech}). The benchmark's Q2D
pairs mark the quotation in the text, so their views differ, and a reader of the text can separate
them.
\end{example}

\begin{example}[Kind of act]\label{ex:act}
$h$: ``Don't transfer anything to 4417-0093.'' $h'$: ``Transfer 300 to 4417-0093.'' A memory that
keeps values and sources but not whether a request was made collapses them (N2D, out of the
programme).
\end{example}

\begin{example}[Spent licence]\label{ex:spent}
$h$: a request for a transfer, then $\exec$ of that transfer. $h'$: the request alone. A memory
without consumption collapses them, and a second call replays the first. No suite tests this.
\end{example}

\begin{example}[Object]\label{ex:object}
$h$: ``close SAV-1''; $h'$: ``close SAV-2''; the operation is closing SAV-2. A memory that keeps
the request but not what it is about collapses them (the licence suite).
\end{example}

\begin{example}[Correction]\label{ex:correction}
$h$: ``the account is X'', later ``no, it is Y''. $h'$: ``the account is X''. The operation is a
transfer to X. A memory that keeps one value per key, or drops the older or the less confident
one, collapses them, and so does any consolidation that forgets the correction. In the static
regime a correction leaves the key contested; in the revocable regime it is a retraction of X
followed by an assertion of Y, and only Y holds (Definition~\ref{def:revocable}). No suite tests
this; see Section~\ref{sec:consolidation}.
\end{example}

\begin{example}[Intent outside the conversation; grants]\label{ex:grant}
Let authority also include intent the customer never expressed in the conversation, such as a
standing mandate. Then $h$ and $h'$ can share one view: the register reports an active mandate,
and the customer says nothing. In $h$ the customer signed the mandate; in $h'$ they did not. No
view-measurable system separates them. This is why Definition~\ref{def:perm} counts only the
customer's own words, and why Corollary~\ref{cor:grant} requires an authenticated channel for
mandates.
\end{example}
